\section{Positive domain of the elliptic chart}
\begin{proposition}[Membership of the representations in the elliptic chart]
\label{iv:ellipse:pertenencia}
The positive branch of \eqref{eq:ii-par-angular}, with the
intervals \eqref{md:accion:intervalos}, satisfies
\(A>0\), \(0<C^*<A\), and \(\hbar_s>0\).
Reversal of the odd coordinate preserves \(|C^*/A|<1\).
\end{proposition}
\begin{proof}
Proposition \ref{md:accion:positividad} gives
\(0<\eta_{\mathrm{ret}}<H_5\) and positivity of both
action sections. For an upper bound on \(H_5\),
\eqref{md:accion:intervalos} implies
\(\alpha<1/125\) and \(\varphi>8/5\). Retaining only
the positive terms of \eqref{eq:el-h5} yields
\[
 H_5<
 \frac{\sqrt5}{2}+
 \frac{9}{16\cdot125^2}+
 \frac{7}{48\cdot125^4}.
\]
Since \(5<81/16\), one has \(\sqrt5/2<9/8\). Moreover,
\[
 \frac{9}{16\cdot125^2}+
 \frac{7}{48\cdot125^4}
 =\frac{421882}{11718750000}<\frac1{1000}.
\]
Thus \(H_5<9/8+1/1000<2\). Using
\(\alpha>7/1000\) now gives
\[
 0<C^*=2(\eta_{\mathrm{ret}}+\alpha)
 <2\left(2+\frac1{125}\right)=\frac{502}{125}
 <7<1000\alpha=A.
\]
The channel interchange in
\ref{prop:ii-inversion-canales-angulares} replaces \(C^*\) by \(-C^*\)
and preserves \(A\); it therefore preserves the absolute inequality.
\end{proof}
